\chapter{Set-Theoretic Foundations}
\label{app:foundations}

Category theory, like all of mathematics, requires some set-theoretic foundation. The main text uses set theory informally, appealing to intuition and the reader's common sense. This appendix discusses the foundational issues that arise when categories become ``large,'' and presents two standard resolutions: the von Neumann-Bernays-Gödel (NBG) approach using classes, and Grothendieck's universe approach.

\section{Naive Set Theory and Its Limits}

In naive set theory, one works with sets freely, forming them by comprehension: $\{x \mid \phi(x)\}$ for any property $\phi$. Russell's paradox (1901) shows this is inconsistent: let $R = \{x \mid x \notin x\}$; then $R \in R \iff R \notin R$.

Zermelo-Fraenkel set theory with Choice (ZFC) avoids this paradox by restricting comprehension and laying down explicit axioms.

\begin{definition}[The axioms of ZFC (informal)]
\leavevmode
\begin{enumerate}
  \item \textbf{Extensionality.} Sets with the same elements are equal: $\forall z(z \in x \iff z \in y) \Rightarrow x = y$.
  \item \textbf{Empty set.} There exists a set with no elements: $\exists x \forall y (y \notin x)$.
  \item \textbf{Pairing.} For any $a, b$, the set $\{a,b\}$ exists.
  \item \textbf{Union.} For any set $\mathcal{A}$, the union $\bigcup \mathcal{A}$ exists.
  \item \textbf{Power set.} For any set $x$, the power set $\mathcal{P}(x) = \{y \mid y \subseteq x\}$ exists.
  \item \textbf{Separation (Restricted Comprehension).} For any set $x$ and formula $\phi$, $\{y \in x \mid \phi(y)\}$ exists.
  \item \textbf{Replacement.} If $\phi(x,y)$ defines a function on a set, the image is a set.
  \item \textbf{Infinity.} There exists an infinite set (containing $\emptyset$ and closed under $x \mapsto x \cup \{x\}$).
  \item \textbf{Foundation (Regularity).} Every nonempty set $x$ has an element disjoint from $x$. (No infinite descending $\in$-chains.)
  \item \textbf{Choice (AC).} Every family of nonempty sets has a choice function.
\end{enumerate}
\end{definition}

In ZFC, all mathematical objects are encoded as sets. The natural numbers are $0 = \emptyset$, $1 = \{\emptyset\}$, $2 = \{\emptyset, \{\emptyset\}\}$, etc.\ (von Neumann ordinals). Ordered pairs are Kuratowski pairs: $(a,b) = \{\{a\}, \{a,b\}\}$. Functions are sets of ordered pairs.

\section{The Size Problem in Category Theory}

A fundamental problem arises: the collection of all sets is not a set (Russell's paradox applies). Therefore $\Set$ does not have an object set—it is a \emph{large} category. Similarly, $\Grp$, $\Top$, etc.\ are large.

But the main concern for category theory is: can we form the collection of all small categories? Can we form the functor category $[\cat{C}\op, \Set]$ when $\cat{C}$ is small? Can we compose large functors?

Two standard approaches resolve this.

\section{Classes: The NBG and MK Approach}

\textbf{von Neumann-Bernays-Gödel set theory (NBG)} and \textbf{Morse-Kelley set theory (MK)} extend ZFC by a primitive notion of \textbf{class}: a class is any collection, and a class that is itself an element of another class is called a \textbf{set}. Classes that are not sets are \textbf{proper classes}.

\begin{itemize}
  \item In NBG/MK, $V$ (the class of all sets) and $\mathbf{On}$ (the class of all ordinals) are proper classes.
  \item The ``class of all groups'' $\text{ob}(\Grp)$ is a proper class; the hom-set $\Grp(G,H)$ is a set.
  \item A \textbf{large category} has a proper class of objects; a \textbf{small category} has a set of objects (and a set of morphisms). A \textbf{locally small} category has a set of morphisms between any two objects.
\end{itemize}

The drawback: in NBG, you cannot form the ``class of all classes'' (there is no class of all large categories). This means $\mathbf{CAT}$ (the category of all large categories) is not well-defined in NBG.

\section{Universes: The Grothendieck Approach}

Grothendieck introduced a cleaner solution using \textbf{universes}.

\begin{definition}[Grothendieck universe]
A \textbf{Grothendieck universe} is a set $\mathcal{U}$ satisfying:
\begin{enumerate}
  \item $x \in \mathcal{U}$ and $y \in x$ implies $y \in \mathcal{U}$ (transitivity).
  \item $x, y \in \mathcal{U}$ implies $\{x,y\} \in \mathcal{U}$ (pairing).
  \item $x \in \mathcal{U}$ implies $\mathcal{P}(x) \in \mathcal{U}$ (power set).
  \item If $I \in \mathcal{U}$ and $\{x_i\}_{i \in I}$ is a family with all $x_i \in \mathcal{U}$, then $\bigcup_i x_i \in \mathcal{U}$ (union).
  \item $\omega \in \mathcal{U}$ (the natural numbers are in $\mathcal{U}$).
\end{enumerate}
A set $x$ is \textbf{$\mathcal{U}$-small} if $x \in \mathcal{U}$.
\end{definition}

\begin{remark}
A Grothendieck universe is essentially a model of ZFC inside ZFC. The existence of a Grothendieck universe is equivalent (over ZFC) to the existence of an \textbf{inaccessible cardinal}: an uncountable cardinal $\kappa$ that is a limit cardinal and regular ($\text{cf}(\kappa) = \kappa$). The existence of inaccessible cardinals cannot be proved in ZFC (by Gödel's incompleteness theorems).
\end{remark}

\textbf{The universe axiom} (Grothendieck's convention) asserts: every set is contained in some Grothendieck universe.

With this axiom (or with a fixed universe $\mathcal{U}_0 \in \mathcal{U}_1 \in \mathcal{U}_2 \in \cdots$), we can:
\begin{itemize}
  \item Call a category \textbf{small} if its objects and morphisms are $\mathcal{U}_0$-small.
  \item Call a category \textbf{large} (or \textbf{locally small}) if its hom-sets are $\mathcal{U}_0$-small.
  \item Form the ``category of all small categories'' $\Cat$ (which is $\mathcal{U}_1$-small).
  \item Form functor categories $[\cat{C}\op, \Set]$ for small $\cat{C}$, which are locally small.
\end{itemize}

Grothendieck's universe approach is now standard in algebraic geometry and topos theory.

\section{The Axiom of Choice}

The Axiom of Choice (AC) is used extensively in category theory:
\begin{itemize}
  \item Choosing skeleta of categories (one representative from each isomorphism class).
  \item Showing that a fully faithful essentially surjective functor is an equivalence (requires choosing the equivalence inverse).
  \item The General Adjoint Functor Theorem (GAFT).
  \item The existence of enough injectives and projectives.
\end{itemize}

The connection between AC and category theory is deep: many categorical theorems are equivalent to AC or weaker choice principles. In constructive mathematics and intuitionistic logic (as in topos theory), AC is not assumed, and many results require reformulation.

\section{Cardinal Arithmetic}

We collect some basic facts about cardinals used in the text.

\begin{definition}
A \textbf{cardinal} is an ordinal that cannot be put in bijection with any smaller ordinal. For a set $X$, $|X|$ denotes the cardinality (the unique cardinal in bijection with $X$, assuming AC).

\begin{itemize}
  \item $\aleph_0 = |\mathbb{N}|$ (countably infinite).
  \item $\aleph_1$ = the first uncountable cardinal.
  \item $\mathfrak{c} = 2^{\aleph_0} = |\mathbb{R}|$ (the continuum).
\end{itemize}

A set is \textbf{finite} if $|X| < \aleph_0$, \textbf{countable} if $|X| \leq \aleph_0$, \textbf{uncountable} otherwise.
\end{definition}

\begin{theorem}[Cantor]
$|X| < |\mathcal{P}(X)|$ for all sets $X$. In particular, there is no ``largest'' cardinal.
\end{theorem}

\begin{definition}
A cardinal $\kappa$ is \textbf{regular} if it cannot be expressed as a sum $\sum_{i < \lambda} \kappa_i$ with $\lambda < \kappa$ and each $\kappa_i < \kappa$. $\aleph_0$ and every successor cardinal $\aleph_{\alpha+1}$ are regular; limit cardinals may or may not be.

A cardinal $\kappa > \aleph_0$ is \textbf{(strongly) inaccessible} if it is regular and $2^\lambda < \kappa$ for all $\lambda < \kappa$.
\end{definition}

Every Grothendieck universe $\mathcal{U}$ corresponds to an inaccessible cardinal $\kappa = |\mathcal{U}|$, and $\mathcal{U} = V_\kappa$ (the $\kappa$-th stage of the cumulative hierarchy).

\section{Summary of Conventions}

Throughout this book, we use the following conventions:
\begin{itemize}
  \item \textbf{Small category}: objects and morphisms form sets.
  \item \textbf{Locally small category}: each hom-class is a set.
  \item \textbf{Large category}: objects may form a proper class (or a set outside $\mathcal{U}_0$).
  \item We assume a suitable universe axiom so that all functor categories we form are legitimate.
  \item We use the Axiom of Choice without explicit mention.
  \item Set-theoretic subtleties beyond local smallness are noted but not belabored.
\end{itemize}
